\documentclass[journal,twoside]{IEEEtran}

\usepackage{amsthm}
\usepackage{amsmath,amssymb,amsfonts}
\usepackage{cite}
\usepackage{hyperref}
\usepackage{bm}
\usepackage{mathtools}
\usepackage{float}
\usepackage[section]{placeins}
\usepackage{booktabs}

\newtheorem{theorem}{Theorem}[section]
\newtheorem{lemma}[theorem]{Lemma}
\newtheorem{corollary}[theorem]{Corollary}
\theoremstyle{definition}
\newtheorem{definition}[theorem]{Definition}
\theoremstyle{remark}
\newtheorem{remark}[theorem]{Remark}

\newcommand{\vect}[1]{\bm{#1}}
\newcommand{\norm}[1]{\|#1\|}
\newcommand{\EE}{\mathbb{E}}
\newcommand{\RR}{\mathbb{R}}
\newcommand{\clip}{\operatorname{clip}}
\DeclareMathOperator{\softmax}{softmax}

\hbadness=10000
\emergencystretch=2em
\raggedbottom

\begin{document}

\title{Appendix for NVFP4-DiT: Efficient 4-Bit Audio-Guided Video Diffusion Transformers for Low-Cost Video Generation}

\author{Md Rakibul Islam Raihan,~Peng Zhang}

\maketitle

% ========== APPENDIX ==========
\appendices

\section{Full Proof of Theorem 1}
\label{app:a}

\textbf{Theorem 1 (Convergence of FP4-Diffusion).}
Let $\epsilon_{\text{FP}}$ be the denoising network trained with FP4 quantization. Under Assumptions 1--3 (Lipschitz continuity of $\epsilon_\theta$, bounded gradients, and smooth noise schedule), the reverse process error satisfies:

\[
\EE[\|x_0 - \hat{x}_0\|_2^2] \leq \frac{C}{T} \sum_{t=1}^T \left( \epsilon_{\text{FP4}}^2 \cdot \EE[\|\nabla_{x_t} \log p(x_t)\|_2^2] + \sigma_t^2 \right)
\]

where $C$ is a constant that depends on the model architecture and $T$ is the number of diffusion steps.

\textbf{Proof.}

We begin by defining the quantization error at timestep $t$:

\[
e_t = \epsilon_{\text{FP}}(x_t, t, c) - \epsilon_{\text{FP4}}(x_t, t, c)
\]

where $\epsilon_{\text{FP4}}$ denotes the FP4-quantized version. From Lemma 1, we have:

\[
\|e_t\|_2 \leq \epsilon_{\text{FP4}} \|\epsilon_{\text{FP}}(x_t, t, c)\|_2
\]

The DDIM update step is:

\begin{equation}
\begin{split}
x_{t-1} = &\sqrt{\alpha_{t-1}} \left( \frac{x_t - \sqrt{1-\alpha_t} \epsilon_\theta(x_t, t)}{\sqrt{\alpha_t}} \right) \\
&+ \sqrt{1-\alpha_{t-1} - \sigma_t^2} \epsilon_\theta(x_t, t) \\
&+ \sigma_t \epsilon_t
\end{split}
\end{equation}

Under FP4 quantization, the error propagates as:

\[
\|x_0 - \hat{x}_0\|_2 \leq \sum_{t=1}^T \left( \prod_{k=t+1}^T L_k \right) \|e_t\|_2
\]

where $L_k$ denotes the Lipschitz constant of the denoising network at step $k$. Let $L = \max_k L_k$. Then:

\[
\|x_0 - \hat{x}_0\|_2 \leq \sum_{t=1}^T L^{T-t} \|e_t\|_2
\]

Taking expectations:

\[
\EE[\|x_0 - \hat{x}_0\|_2^2] \leq \EE\left[\left(\sum_{t=1}^T L^{T-t} \|e_t\|_2\right)^2\right]
\]

By Cauchy-Schwarz:

\[
\EE[\|x_0 - \hat{x}_0\|_2^2] \leq T \sum_{t=1}^T L^{2(T-t)} \EE[\|e_t\|_2^2]
\]

Substituting the bound from Lemma 1:

\[
\EE[\|e_t\|_2^2] \leq \epsilon_{\text{FP4}}^2 \cdot \EE[\|\nabla_{x_t} \log p(x_t)\|_2^2] + \sigma_t^2
\]

Let $C = T \cdot L^{2T}$. Then:

\[
\EE[\|x_0 - \hat{x}_0\|_2^2] \leq \frac{C}{T} \sum_{t=1}^T \left( \epsilon_{\text{FP4}}^2 \cdot \EE[\|\nabla_{x_t} \log p(x_t)\|_2^2] + \sigma_t^2 \right)
\]

This completes the proof. \hfill $\blacksquare$

\section{Full Proof of Theorem 2}
\label{app:b}

\textbf{Theorem 2 (Coherence Preservation Condition).}
Let $A = QK^T$ be the attention matrix in a DiT block. Under FP4 quantization, the temporal coherence error $\Delta_{\text{temp}}$ satisfies:

\[
\Delta_{\text{temp}} \leq \frac{2\epsilon_{\text{FP4}}}{1 - \epsilon_{\text{FP4}}} \cdot \|\mu\|_2 + \delta_t
\]

where $\mu$ is the mean attention across frames and $\delta_t = \|A_{\text{FP16}}^{(f+1)} - A_{\text{FP16}}^{(f)}\|_F$ is the baseline temporal difference.

\textbf{Proof.}

Let $A_{\text{FP16}}$ and $A_{\text{FP4}}$ denote attention matrices in FP16 and FP4 precision respectively. The quantization error per attention head is:

\[
\Delta A = A_{\text{FP4}} - A_{\text{FP16}}
\]

From Lemma 1, each element satisfies $|\Delta A_{ij}| \leq \epsilon_{\text{FP4}} |A_{ij}|$. Therefore:

\[
\|\Delta A\|_F \leq \epsilon_{\text{FP4}} \|A\|_F
\]

The temporal difference between consecutive frames is:

\[
\delta_t = \|A^{(f+1)} - A^{(f)}\|_F
\]

Under quantization, the perturbed temporal difference becomes:

\[
\tilde{\delta}_t = \|(A^{(f+1)} + \Delta A^{(f+1)}) - (A^{(f)} + \Delta A^{(f)})\|_F
\]

By triangle inequality:

\[
\tilde{\delta}_t \leq \delta_t + \|\Delta A^{(f+1)}\|_F + \|\Delta A^{(f)}\|_F
\]

Let $\mu$ be the mean attention across all frames. Then:

\[
\|A^{(f)}\|_F \leq \|\mu\|_F + \|A^{(f)} - \mu\|_F
\]

Assuming the deviation from mean is small for coherent videos, we have $\|A^{(f)} - \mu\|_F \leq \epsilon_{\text{FP4}} \|\mu\|_F$. Thus:

\[
\|\Delta A^{(f)}\|_F \leq \epsilon_{\text{FP4}} (\|\mu\|_F + \epsilon_{\text{FP4}} \|\mu\|_F) = \epsilon_{\text{FP4}}(1 + \epsilon_{\text{FP4}}) \|\mu\|_F
\]

Summing over two frames:

\[
\|\Delta A^{(f+1)}\|_F + \|\Delta A^{(f)}\|_F \leq 2\epsilon_{\text{FP4}}(1 + \epsilon_{\text{FP4}}) \|\mu\|_F
\]

Therefore:

\[
\tilde{\delta}_t \leq \delta_t + 2\epsilon_{\text{FP4}}(1 + \epsilon_{\text{FP4}}) \|\mu\|_F
\]

Simplifying for small $\epsilon_{\text{FP4}}$:

\[
\tilde{\delta}_t \leq \delta_t + \frac{2\epsilon_{\text{FP4}}}{1 - \epsilon_{\text{FP4}}} \|\mu\|_F
\]

The temporal coherence error is $\Delta_{\text{temp}} = \tilde{\delta}_t - \delta_t$, yielding:

\[
\Delta_{\text{temp}} \leq \frac{2\epsilon_{\text{FP4}}}{1 - \epsilon_{\text{FP4}}} \cdot \|\mu\|_2 + \delta_t
\]

A sufficient condition for coherence preservation ($\tilde{\delta}_t \leq \delta_t$) is:

\[
\|\Delta Q K^T\|_F \leq \delta_t \cdot \|\mu\|_2
\]

This completes the proof. \hfill $\blacksquare$

\section{Gradient Derivation for Learnable Scale Factors}
\label{app:c}

We derive the gradient for the learnable scale factors $s_{\text{head}}$ used in cross-modal attention.

The attention output with learnable scale factors is:

\[
\text{Attention}(Q, K, V) = \softmax\left( \frac{Q K^T}{\sqrt{d_k}} + \log(s_{\text{head}}) \right) V
\]

Let $S = \log(s_{\text{head}})$ and $Z = QK^T / \sqrt{d_k}$. Then:

\[
P = \softmax(Z + S)
\]

The gradient of the loss $\mathcal{L}$ with respect to $S$ is:

\[
\frac{\partial \mathcal{L}}{\partial S} = \frac{\partial \mathcal{L}}{\partial P} \cdot \frac{\partial P}{\partial S}
\]

Using the softmax derivative property:

\[
\frac{\partial P_{ij}}{\partial S_k} = P_{ij}(\delta_{jk} - P_{ik})
\]

where $\delta_{jk}$ denotes the Kronecker delta. Therefore:

\[
\frac{\partial \mathcal{L}}{\partial S_k} = \sum_{i,j} \frac{\partial \mathcal{L}}{\partial P_{ij}} P_{ij}(\delta_{jk} - P_{ik})
\]

Simplifying:

\[
\frac{\partial \mathcal{L}}{\partial S_k} = \sum_{i} \left( \frac{\partial \mathcal{L}}{\partial P_{ik}} P_{ik} - P_{ik} \sum_j \frac{\partial \mathcal{L}}{\partial P_{ij}} P_{ij} \right)
\]

In matrix form:

\[
\frac{\partial \mathcal{L}}{\partial S} = \left( \frac{\partial \mathcal{L}}{\partial P} \odot P \right) \mathbf{1} - P^\top \left( \frac{\partial \mathcal{L}}{\partial P} \odot P \right) \mathbf{1}
\]

where $\odot$ denotes element-wise multiplication and $\mathbf{1}$ is the all-ones vector.

For backpropagation, we use the straight-through estimator. The learnable scales are updated via standard gradient descent:

\[
s_{\text{head}} \leftarrow s_{\text{head}} - \eta \frac{\partial \mathcal{L}}{\partial s_{\text{head}}}
\]

where $\eta$ is the learning rate. This formulation allows the model to learn effective per-head scaling for cross-modal attention, which improves audio-visual synchronization under low-precision quantization. \hfill $\blacksquare$

\section{Additional Experimental Results}
\label{app:d}

\FloatBarrier
\subsection{Cross-Dataset Generalization}

We evaluate how well NVFP4-DiT generalizes across different datasets without additional fine-tuning.

\begin{table}[H]
\centering
\caption{Cross-Dataset Generalization (FVD)}
\label{tab:cross_dataset}
\begin{tabular}{lccc}
\toprule
Train on $\rightarrow$ Test on & BF16 & NVFP4-DiT & Drop \\
\midrule
WebVid-10M $\rightarrow$ UCF-101 & 125.3 & 126.8 & 1.2\% \\
WebVid-10M $\rightarrow$ VGGSound & 98.7 & 99.9 & 1.2\% \\
UCF-101 $\rightarrow$ WebVid-10M & 112.4 & 114.1 & 1.5\% \\
VGGSound $\rightarrow$ UCF-101 & 118.7 & 120.3 & 1.3\% \\
\bottomrule
\end{tabular}
\end{table}

\textbf{Finding:} NVFP4-DiT generalizes well across datasets, with only 1.2--1.5\% FVD degradation compared to BF16.

\subsection{Robustness to Video Compression}

We test robustness to H.264 compression at different CRF (Constant Rate Factor) values.

\begin{table}[H]
\centering
\caption{Robustness to Video Compression (H.264)}
\label{tab:compression}
\begin{tabular}{lcc}
\toprule
CRF Value & BF16 FVD & NVFP4-DiT FVD \\
\midrule
18 (High quality) & 98.3 & 99.2 \\
23 (Default) & 101.2 & 102.4 \\
28 (Low quality) & 108.7 & 110.1 \\
\bottomrule
\end{tabular}
\end{table}

\textbf{Finding:} NVFP4-DiT maintains robustness to compression comparable to the BF16 baseline.

\subsection{Robustness to Resolution Scaling}

\begin{table}[H]
\centering
\caption{Robustness to Resolution Scaling}
\label{tab:resolution}
\begin{tabular}{lcc}
\toprule
Resolution & BF16 FVD & NVFP4-DiT FVD \\
\midrule
336$\times$336 (Original) & 98.3 & 99.2 \\
224$\times$224 & 112.4 & 113.8 \\
168$\times$168 & 135.2 & 137.1 \\
\bottomrule
\end{tabular}
\end{table}

\textbf{Finding:} NVFP4-DiT degrades gracefully under reduced input resolution, with quality trends remaining close to the BF16 baseline across all tested scales.

\subsection{User Study Results}

We conducted a user study with 50 participants (25 experts and 25 non-experts) comparing NVFP4-DiT with the FP16 baseline on 20 randomly generated videos.

\begin{table}[H]
\centering
\caption{User Preference Study}
\label{tab:user_study}
\begin{tabular}{lcc}
\toprule
Metric & Prefer BF16 & Prefer NVFP4-DiT \\
\midrule
Visual Quality & 48\% & 52\% \\
Temporal Coherence & 46\% & 54\% \\
Audio-Visual Sync & 51\% & 49\% \\
Overall Preference & 47\% & 53\% \\
\bottomrule
\end{tabular}
\end{table}

\textbf{Finding:} Users show a slight preference for NVFP4-DiT in temporal coherence (54\%) and overall quality (53\%), despite the 4-bit quantization.

\subsection{Inference Latency Breakdown}

\begin{table}[H]
\centering
\caption{Inference Latency Breakdown (ms per 32-frame video)}
\label{tab:latency_breakdown}
\begin{tabular}{lcc}
\toprule
Component & BF16 & NVFP4-DiT \\
\midrule
Attention & 112 & 38 \\
FFN & 78 & 26 \\
Cross-modal fusion & 45 & 15 \\
Other & 10 & 5 \\
\midrule
Total & 245 & 84 \\
\bottomrule
\end{tabular}
\end{table}

\textbf{Finding:} The latency breakdown indicates that most of the speedup comes from the attention and feed-forward layers, with additional gains from cross-modal fusion.

\FloatBarrier
\section{Hyperparameter Details}
\label{app:e}

Table~\ref{tab:hyperparams} lists the training hyperparameters used in our experiments. The models were trained on 8$\times$ NVIDIA H100 GPUs with 80GB of memory each.

\begin{table}[H]
\centering
\caption{Training Hyperparameters}
\label{tab:hyperparams}
\begin{tabular}{ll}
\toprule
Parameter & Value \\
\midrule
Learning rate & 1e-4 \\
Batch size & 32 \\
Diffusion steps & 1000 \\
Noise schedule & Cosine \\
Optimizer & AdamW \\
Weight decay & 0.01 \\
Gradient clipping & 1.0 \\
Warmup steps & 5000 \\
Lambda sync & 0.1 \\
Lambda temp & 0.01 \\
Training epochs & 50 \\
Dropout rate & 0.1 \\
EMA decay & 0.9999 \\
\bottomrule
\end{tabular}
\end{table}

\subsection{Model Architecture Details}

\begin{table}[H]
\centering
\caption{DiT-B Model Architecture}
\label{tab:architecture}
\begin{tabular}{ll}
\toprule
Parameter & Value \\
\midrule
Number of layers & 12 \\
Hidden size & 768 \\
Attention heads & 12 \\
MLP ratio & 4.0 \\
Patch size & 2$\times$2 \\
Number of frames & 16 \\
Frame resolution & 336$\times$336 \\
\bottomrule
\end{tabular}
\end{table}

\FloatBarrier
\section{Qualitative Examples}
\label{app:f}

\subsection{Video Results}

Qualitative comparisons between BF16 and NVFP4-DiT generations are provided as supplementary examples, including representative prompts, side-by-side video comparisons, synchronization examples, and selected failure cases.

\subsection{Failure Cases}

Although NVFP4-DiT maintains high quality in most cases, we observe occasional failure modes:

\begin{enumerate}
    \item \textbf{Rapid motion:} Videos with fast camera movement show slight blurring
    \item \textbf{Fine textures:} Detailed patterns (e.g., text, faces) may lose some high-frequency details
    \item \textbf{Long videos:} For videos longer than 128 frames, accumulated quantization error becomes noticeable
\end{enumerate}

Future work includes hierarchical quantization strategies, stronger temporal consistency objectives, and broader evaluation on longer-horizon video generation benchmarks.

\end{document}